\chapter{Contexto y motivación}

\begin{quotation}[Físico teórico]{Richard P. Feynman (1918--1988)}
Lo que no puedo crear, no lo entiendo.
\end{quotation}

\begin{abstract}
Antes de construir un edificio, es necesario conocer sus cimientos. Este capítulo presenta el contexto sobre el que se asienta el proyecto.
\end{abstract}

\section{El mundo de la robótica móvil}

La robótica móvil estudia el diseño y control de sistemas que pueden desplazarse de forma autónoma en un entorno físico. A diferencia de otros robots que trabajan en espacios muy restringidos o guiados por carriles, un robot móvil se mueve en un espacio continuo, donde el entorno puede cambiar, y donde ningún trayecto es exactamente igual al anterior. Esto hace que la navegación autónoma sea un problema desafiante: se trata no solo de moverse, sino de moverse de forma segura, eficiente y con un objetivo claro.

En este contexto, los robots móviles se utilizan en entornos muy diversos: fábricas, almacenes, hospitales, museos, edificios universitarios o incluso en la calle. En todos ellos, el robot debe ser capaz de interpretar su entorno, saber dónde se encuentra, y planificar rutas hacia lugares de destino sin chocar ni bloquear el paso a personas u otros sistemas. La clave no está solo en tener ruedas, sensores o motores, sino en integrarlos en un sistema que tome decisiones razonables sobre el movimiento.

\section{Navegación autónoma: entre la percepción y la decisión}

La navegación autónoma de un robot móvil puede resumirse como un ciclo continuo de tres pasos: percibir el entorno, localizarse en él y decidir hacia dónde moverse. En este proyecto, el robot se apoya en sensores como el láser o el SLAM para construir un mapa 2D del entorno, y utiliza algoritmos de planificación de trayectoria para calcular rutas desde el punto actual hasta un objetivo. En este sentido, como señaló Alfred Korzybski, ``\emph{el mapa no es el territorio}''\cite{korzybski1933}, por lo que esta representación nunca es perfecta y el sistema debe ser capaz de adaptarse a las diferencias entre el modelo y la realidad. La decisión final se traduce en comandos de velocidad que el robot ejecuta, ajustando su trayectoria en tiempo real.

Esta idea se acerca mucho a la frase que abre el capítulo: la verdadera medida de la autonomía no es que el robot pueda moverse, sino que sepa \emph{hacia dónde ir} y cómo llegar de forma segura. Un sistema de navegación autónoma no es un mero motor que avanza recto, sino un sistema de control que reacciona a obstáculos, cambios de planificación y errores de localización, ajustando su comportamiento en cada instante.

\section{ROS 2 como base del sistema}

ROS 2 (Robot Operating System 2) es un marco de desarrollo de software ampliamente utilizado en robótica por su arquitectura distribuida, modular y orientada a componentes. En lugar de construir un sistema monolítico, ROS 2 permite que cada parte del robot (percepción, localización, planificación, control y supervisión) se implemente como un nodo independiente que se comunica con el resto mediante mensajes. Esto favorece la reutilización de código, la integración de sensores y la extensión futura del sistema.

En este trabajo se utiliza la distribución \textbf{ROS 2 Humble}, que ofrece estabilidad y soporte de largo plazo, ideal para proyectos académicos y experimentales. Además, ROS 2 proporciona herramientas estándar para la depuración, la visualización y el testeo de nodos, lo que permite entender el comportamiento del sistema en tiempo de ejecución y corregir fallos de forma más eficiente. Desde ROS 2 se coordinan tanto el robot simulado en Ignition Gazebo Fortress (el simulador por antonomasia del ecosistema ROS2 Humble) como el sistema real de TurtleBot 4, permitiendo que el desarrollo se transfiera con relativa facilidad entre simulación y hardware.

\section{TurtleBot 4: plataforma para navegar y validar}

TurtleBot 4 es una plataforma de robótica móvil diseñada para investigación, aprendizaje y prototipado. Combina un chasis móvil, sensores de proximidad, láser (o equivalentes), procesamiento embebido y conectividad inalámbrica, y está pensada precisamente para experimentar con navegación autónoma, mapeado y control de robots. Su compatibilidad con ROS 2 Humble facilita la integración de paquetes de navegación, SLAM y visualización, lo que la convierte en una base muy adecuada para un TFG centrado en navegación autónoma.

En este proyecto, TurtleBot 4 se utiliza como laboratorio rodante de manera simulada, dejando las pruebas en el robot físico como trabajo futuro en base a lo aprendido. El entorno de simulación permite probar algoritmos, parámetros y configuraciones de forma segura, mientras que el robot real permitiría validar que el comportamiento en el mundo físico es coherente con el esperado en simulación.

\section{Del mapeado a la navegación continua}

Un paso clave en la navegación autónoma es la construcción de un mapa del entorno. Para ello se utiliza \textbf{SLAM} (Simultaneous Localization And Mapping), que permite que el robot se localice y construya un mapa 2D a la vez que se mueve. Una vez existe un mapa fiable, el robot puede usarlo para planificar trayectorias desde su posición actual hasta un objetivo, evitando obstáculos y respetando el espacio visible. Además, el sistema de navegación puede reaccionar ante cambios, como nuevos obstáculos o cambios de planificación, recalculando la ruta en tiempo real.

En este contexto, la navegación no es un evento único, sino un proceso continuo. El robot debe ser capaz de replanificar, recuperarse de situaciones de bloqueo, respetar zonas de seguridad y adaptar su comportamiento al entorno. Esto implica que el diseño del sistema no solo tenga un “modo de avance”, sino estrategias de recuperación, de detección de problemas y de replanificación. En este trabajo, la pila de navegación \textbf{Nav2} proporciona este conjunto de herramientas dentro de ROS 2, permitiendo que el robot navegue de forma autónoma en entornos estáticos y dinámicos.

\section{Estado del arte}

Actualmente, la navegación autónoma de robots móviles ha dejado de ser una tecnología experimental para convertirse en una realidad en entornos de fábrica, logística y edificios inteligentes. Muchas empresas y proyectos de investigación utilizan sistemas basados en ROS 2, Nav2 y SLAM para dotar a sus robots de una navegación robusta y flexible. A su vez, plataformas como TurtleBot 4 ofrecen una forma accesible de experimentar con estos sistemas, lo que las convierte en candidatas naturales para proyectos académicos y de pruebas.

En el ámbito académico, se han desarrollado múltiples trabajos que analizan el rendimiento de distintos algoritmos de navegación, modos de mapeado y estrategias de control para robots móviles. Estos trabajos muestran que la elección de la pila de navegación, los parámetros de planificación y la calidad del mapa tienen un impacto directo en la eficiencia y seguridad del robot. Además, la combinación de simulación y hardware real se ha convertido en una práctica estándar, ya que permite probar configuraciones complejas sin riesgo de daño físico.

En este contexto, el proyecto se inserta en la línea de trabajo que explora la navegación autónoma de robots móviles con ROS 2 y TurtleBot 4, centrándose en la implementación, evaluación y ajuste de algoritmos de navegación y mapeado. La hipótesis central es que, mediante la configuración adecuada de la pila de navegación y la integración de SLAM, se puede obtener un sistema de navegación autónoma robusto, capaz de desplazarse de forma segura en entornos de interior, tanto en simulación como en el mundo real.
